\title{MInDes 软件开发入门手册}
\author{
    王梓航\thanks{主笔, E-mail: \bhref{mailto:w.zihang@qq.com}{w.zihang@qq.com}}
    \and 黄奇\thanks{审阅, E-mail: \bhref{mailto:qihuang0908@163.com}{qihuang0908@163.com}}
}
\date{\today}
\maketitle
\tableofcontents
\newpage

\section{简介}\label{Intro}
MInDes (Material Intelligent Design) 程序的开发以运行高效、结构清晰、功能完善为目标, 因涉及第三方库以及多平台开发
等原因, 对软件的开发环境有一定的要求. 本文将介绍 Windows平台的开发环境搭建, Linux平台
开发环境搭建, Windows与Linux平台的程序调试, 基于Git与GitHub的多人协同等, 并在文末补充介绍一些程序开发的基本概念以及一些
可能常见的问题. 旨在分享必要的程序开发知识, 帮助读者快速上手入门程序开发. 欢迎对本文进行补充, 遇到的疑问也欢迎同笔者交流沟通.

笔者的部分软件并未进行汉化, 涉及具体操作时将贴上英文以及自行翻译的中文以提供对照. 翻译的中文可能有所差错, 敬请谅解. 文中出现的
序号0均意为检查项目, 请检查该项以确保接下来的步骤的顺利完成. 文中使用不同的字体以区分不同的对象, 以下是对照条目.
\begin{itemize}
    \item \emph{斜体(楷体) italic} 指代文件, 运行库. 同时也作为强调之用.
    \item {\bfseries{黑体} boldface} 指代文件夹.
    \item {\sffamily{无衬线体} sans serif} 指代程序中出现的图标文字等. 通常有笔者翻译的中文和英文对照. 其中出现的箭头{\textrightarrow}指代
          目录的层级结构.
    \item \cverb|灰底衬线体 code| 指代命令行中输入的命令或是按键快捷键.
\end{itemize}
截至本手册成稿时, 笔者并没有新设备来一一检验文中所述均为无错漏或多余的必要步骤, 如有何处错漏, 敬请批评指正.
如需要快速搭建好开发测试环境, 请移步至章节\ref{QA}: Q\&A 部分查阅您需要的内容.

\endinput